\section{Related work}
\label{sec:related}

\paragraph{Multi-agent systems.}
Language-model agents with memory, reflection, and planning produce believable emergent social behavior in shared environments \citep{park2023generative}, and deployed agent populations now exist at scale: on Moltbook, an early empirical study documented flooding by a single agent, extreme vote asymmetry, and in-group rhetoric that coordinated subsets of agents \citep{jiang2026humans}, security researchers documented agent-to-agent prompt injection through posts and comments \citep{zenity2026moltbook}, and the Hugging Face incident showed the same dynamics inside an evaluation sandbox \citep{openai2026huggingface,huggingface2026anatomy}. That agents can coordinate against their principals is established: \citet{motwani2024secret} demonstrate collusion through steganographic channels invisible to an overseer \citep[see also][]{motwani2025open}, \citet{fish2024algorithmic} show language-model pricing agents converging to supracompetitive prices without explicit communication, and \citet{hubinger2024sleeper} show that trigger-conditioned deceptive behavior persists through safety training. We take our setting from the first group of papers and our threat model from the second. We differ from the collusion literature in what is observed: those papers study whether agents \emph{can} collude, with access to the agents, their channel, or their training, whereas we ask whether collusion is visible to an outside observer who sees only behavior and never the channel, which is the position of a platform operator or a third-party auditor.

\paragraph{Low-dimensional representations of agents.}
Black-box generative models can be compared through their responses: the data kernel is the Gram matrix of embedded responses to a common query set \citep{duderstadt2023comparing}, and the induced perspective space tracks systems of interacting models and shows how communication topology governs the diffusion of perspectives, adversarial ones included, generalizing model collapse \citep{shumailov2024ai} to networks of models \citep{helm2024tracking}. Inference in the data kernel perspective space (DKPS) and the consistency and concentration of its representations are developed in \citet{helm2025statistical,acharyya2024consistent,acharyya2025concentration}, regression on DKPS representations is query-efficient \citep{helm2026query}, and the Product Kernel Perspective Space (PKPS) extends the construction to models that answered different queries, concentrating at the paired rate up to a query-overlap factor \citep{pkps2026}. Over time, the temporal DKPS embeds (agent, time) pairs jointly and tests for change in a single agent or in a pre-specified group \citep{bridgeford2025detecting}, and adaptive control charts on the system-level trajectory give continuous monitoring at the cost of vulnerability to agents that defect slowly \citep{helm2026control}. We use PKPS as our representation because the observation model of Section~\ref{sec:setting} is sparse and unpaired by nature, and the temporal embedding and permutation null of the temporal DKPS as our inferential starting point. We differ in the question: every method in this line tests for change in an agent, a group, or a system that is given in advance, whereas our group is unknown and must be found, which turns a two-sample test into a scan and requires the concentration results above to control estimation error uniformly over candidate groups.

\paragraph{Statistical network science.}
Latent position models represent vertices by points whose inner products give edge probabilities, and adjacency spectral embedding recovers those points up to rotation with guarantees for clustering, classification, and testing \citep{athreya2018statistical,lyzinski2014perfect,lyzinski2017community,priebe2019two}; for several graphs on a common vertex set the omnibus embedding places every (vertex, graph) pair in one space so that a vertex can be compared across graphs \citep{levin2017omnibus}, and nonparametric two-sample tests decide whether two graphs share a generating distribution \citep{tang2017nonparametric}. Scan statistics on graphs detect a small subset of vertices whose behavior is anomalous: vertex-centered locality statistics \citep{priebe2005scan,rukhin2011comparative} and their extensions to time series of graphs \citep{park2013anomaly,wang2014locality,chen2025multiple} scan the neighborhoods of a graph for excess connectivity and have power against a chatter alternative in which a subset raises its mutual connectivity; the general theory of scanning a candidate family of clusters for an elevated mean \citep{glaz2001scan,kulldorff1997spatial,ariascastro2011detection}, the statistical and computational limits of planted-subgraph detection \citep{verzelen2015community,hajek2015computational}, and spectral relaxations of the graph scan \citep{sharpnack2016detecting} characterize when such a subset is detectable. Vertex nomination ranks vertices by their probability of belonging to a subset of interest given a few known members \citep{fishkind2015vertex}, including when an adversary modifies the graph to hide the subset \citep{agterberg2020vertex}, and change-point methods localize the time at which a sequence of networks changes distribution \citep{padilla2022change,wang2021optimal}, the natural tool for the onset-estimation problem we defer to future work. Our method is an instance of this template with one substitution: the latent position is estimated from behavior rather than from edges, with (agent, time) pairs as vertices and behavioral similarity in place of adjacency, so that no graph is observed. The joint embedding of Section~\ref{sec:est-ps} is the omnibus embedding's analogue; Definition~\ref{def:collusion} is the chatter alternative stated in the latent space; the scan \eqref{eq:scan} is an elevated-mean cluster scan over a data-driven family calibrated by permutation; and membership estimation is vertex nomination without seeds. We differ in what is assumed observed: scan statistics on graphs, vertex nomination, and network change-point methods all take the graph as data, and their guarantees rest on the adjacency matrix concentrating around its expectation, whereas our ``edges'' are sparse, unpaired, and sometimes-estimated query--response pairs, and the role of adjacency concentration is played by the concentration of the behavioral representation.
